\documentclass[11pt]{article}
\usepackage{lmodern}
% jugeo-common.tex — shared preamble for all Judgment Geometry papers
% Include via: % jugeo-common.tex — shared preamble for all Judgment Geometry papers
% Include via: % jugeo-common.tex — shared preamble for all Judgment Geometry papers
% Include via: \input{jugeo-common}

% ── Packages ────────────────────────────────────────────────────────────
\usepackage{amsmath,amssymb,amsthm,mathtools}
\usepackage{tikz-cd}
\usepackage{listings}
\usepackage{xcolor}
\usepackage{xspace}
\usepackage{booktabs}
\usepackage{multirow}
\usepackage{enumitem}
\usepackage[expansion=false]{microtype}
\usepackage{hyperref}
\usepackage{cleveref}
% \usepackage{thmtools}  % disabled: not available in this TeX installation
\usepackage{float}
\usepackage{caption}
\usepackage{subcaption}
\usepackage{tabularx}
\usepackage{stmaryrd}       % \llbracket, \rrbracket
\usepackage{bbm}            % \mathbbm{1}

% ── Page geometry ───────────────────────────────────────────────────────
\usepackage[margin=1in]{geometry}

% ── Theorem environments ────────────────────────────────────────────────
\theoremstyle{plain}
\newtheorem{theorem}{Theorem}[section]
\newtheorem{lemma}[theorem]{Lemma}
\newtheorem{proposition}[theorem]{Proposition}
\newtheorem{corollary}[theorem]{Corollary}

\theoremstyle{definition}
\newtheorem{definition}[theorem]{Definition}
\newtheorem{example}[theorem]{Example}
\newtheorem{construction}[theorem]{Construction}

\theoremstyle{remark}
\newtheorem{remark}[theorem]{Remark}
\newtheorem{notation}[theorem]{Notation}

% ── Python listings style ──────────────────────────────────────────────
\definecolor{jugeoBlue}{HTML}{2563EB}
\definecolor{jugeoGreen}{HTML}{059669}
\definecolor{jugeoRed}{HTML}{DC2626}
\definecolor{jugeoPurple}{HTML}{7C3AED}
\definecolor{jugeoGray}{HTML}{6B7280}
\definecolor{jugeoBg}{HTML}{F8FAFC}

\lstdefinestyle{jugeo-python}{
  language=Python,
  basicstyle=\ttfamily\small,
  keywordstyle=\color{jugeoBlue}\bfseries,
  stringstyle=\color{jugeoGreen},
  commentstyle=\color{jugeoGray}\itshape,
  numberstyle=\tiny\color{jugeoGray},
  backgroundcolor=\color{jugeoBg},
  numbers=left,
  numbersep=6pt,
  frame=single,
  framerule=0.4pt,
  rulecolor=\color{jugeoGray!40},
  breaklines=true,
  breakatwhitespace=true,
  showstringspaces=false,
  tabsize=4,
  xleftmargin=1.5em,
  framexleftmargin=1.5em,
  morekeywords={dataclass,frozen,field,Enum,IntEnum,auto,Any,
                Mapping,Sequence,Optional,match,case},
  literate={→}{{$\to$}}1 {←}{{$\leftarrow$}}1
           {≤}{{$\leq$}}1 {≥}{{$\geq$}}1 {≠}{{$\neq$}}1
           {∈}{{$\in$}}1 {∀}{{$\forall$}}1 {∃}{{$\exists$}}1
           {⊕}{{$\oplus$}}1 {⊖}{{$\ominus$}}1,
  escapeinside={(*@}{@*)},
}
\lstset{style=jugeo-python}

\lstdefinestyle{jugeo-output}{
  basicstyle=\ttfamily\small,
  backgroundcolor=\color{jugeoBg},
  frame=single,
  framerule=0.4pt,
  rulecolor=\color{jugeoGray!40},
  breaklines=true,
  showstringspaces=false,
  xleftmargin=1.5em,
  framexleftmargin=0.5em,
  numbers=none,
}

% ── JuGeo-specific macros ──────────────────────────────────────────────

% Judgment 8-tuple
\newcommand{\J}{\mathcal{J}}
\newcommand{\Jud}[8]{(#1,\,#2,\,#3,\,#4,\,#5,\,#6,\,#7,\,#8)}
\newcommand{\JudShort}[1]{\J_{#1}}

% Components
\newcommand{\coord}{\mathsf{c}}            % coordinate
\newcommand{\prop}{\varphi}                 % proposition
\newcommand{\carrier}{\mathcal{A}}          % carrier type
\newcommand{\evid}{\mathcal{E}}             % evidence bundle
\newcommand{\oblig}{\mathcal{O}}            % obligations
\newcommand{\obstr}{\mathcal{B}}            % obstructions
\newcommand{\trust}{\mathcal{T}}            % trust annotation
\newcommand{\prov}{\Pi}                     % provenance

% Trust algebra
\newcommand{\Talg}{\mathfrak{T}}            % trust ordered algebra
\newcommand{\Eadm}{\mathcal{E}_{\mathrm{adm}}}
\newcommand{\tleq}{\preceq}                 % trust ordering
\newcommand{\tjoin}{\oplus}                 % conservative join
\newcommand{\tatten}{\ominus}               % attenuation
\newcommand{\tpromote}{\uparrow_\pi}        % promotion
\newcommand{\tdemote}{\downarrow_\chi}       % demotion

% Trust tiers
\newcommand{\Tcontra}{\ensuremath{\mathsf{CONTRADICTED}}}
\newcommand{\Tunver}{\ensuremath{\mathsf{UNVERIFIED}}}
\newcommand{\Tcopilot}{\ensuremath{\mathsf{COPILOT}}}
\newcommand{\Toracle}{\ensuremath{\mathsf{ORACLE}}}
\newcommand{\Truntime}{\ensuremath{\mathsf{RUNTIME}}}
\newcommand{\Tsolver}{\ensuremath{\mathsf{SOLVER}}}
\newcommand{\Tproof}{\ensuremath{\mathsf{PROOF}}}

% Site and geometry
\newcommand{\Site}{\mathcal{S}}             % semantic site
\newcommand{\Cov}{\mathrm{Cov}}            % covering family
\newcommand{\Ob}{\mathrm{Ob}}              % objects
\newcommand{\Mor}{\mathrm{Mor}}            % morphisms
\newcommand{\res}{\mathrm{res}}            % restriction
\newcommand{\Cech}{\check{C}}              % Čech complex
\newcommand{\Hcoh}[1]{H^{#1}}             % cohomology group

% Descent
\newcommand{\descent}{\mathrm{desc}}
\newcommand{\glue}{\mathrm{glue}}
\newcommand{\overlap}[2]{#1 \cap #2}

% Semantic moves
\newcommand{\VERIFY}{\textsc{Verify}}
\newcommand{\CONSTRUCT}{\textsc{Construct}}
\newcommand{\NORMALIZE}{\textsc{Normalize}}
\newcommand{\GLUE}{\textsc{Glue}}
\newcommand{\RESTRICT}{\textsc{Restrict}}
\newcommand{\TRANSPORT}{\textsc{Transport}}
\newcommand{\REPAIR}{\textsc{Repair}}
\newcommand{\PROMOTE}{\textsc{Promote}}
\newcommand{\CHALLENGE}{\textsc{Challenge}}
\newcommand{\ESCALATE}{\textsc{Escalate}}

% Fragments
\newcommand{\QFLIA}{\texttt{QF\_LIA}}
\newcommand{\QFLRA}{\texttt{QF\_LRA}}
\newcommand{\QFBV}{\texttt{QF\_BV}}
\newcommand{\QFUF}{\texttt{QF\_UF}}

% Misc
\newcommand{\Python}{\textsc{Python}}
\newcommand{\JuGeo}{\textsc{JuGeo}}
\newcommand{\LEAN}{\textsc{Lean}\,4}
\newcommand{\Fstar}{F$^{\star}$}
\newcommand{\Coq}{\textsc{Coq}}
\newcommand{\Dafny}{\textsc{Dafny}}

% Common shortcuts
\newcommand{\ie}{i.e.\@\xspace}
\newcommand{\eg}{e.g.\@\xspace}
\newcommand{\cf}{cf.\@\xspace}
\newcommand{\wrt}{w.r.t.\@\xspace}

% Evaluation shortcuts
\newcommand{\numcases}{300}
\newcommand{\numspec}{100}
\newcommand{\numeq}{100}
\newcommand{\numbug}{100}
\newcommand{\medtime}{6.5\,ms}
\newcommand{\totaltime}{1.949\,s}


% ── Packages ────────────────────────────────────────────────────────────
\usepackage{amsmath,amssymb,amsthm,mathtools}
\usepackage{tikz-cd}
\usepackage{listings}
\usepackage{xcolor}
\usepackage{xspace}
\usepackage{booktabs}
\usepackage{multirow}
\usepackage{enumitem}
\usepackage[expansion=false]{microtype}
\usepackage{hyperref}
\usepackage{cleveref}
% \usepackage{thmtools}  % disabled: not available in this TeX installation
\usepackage{float}
\usepackage{caption}
\usepackage{subcaption}
\usepackage{tabularx}
\usepackage{stmaryrd}       % \llbracket, \rrbracket
\usepackage{bbm}            % \mathbbm{1}

% ── Page geometry ───────────────────────────────────────────────────────
\usepackage[margin=1in]{geometry}

% ── Theorem environments ────────────────────────────────────────────────
\theoremstyle{plain}
\newtheorem{theorem}{Theorem}[section]
\newtheorem{lemma}[theorem]{Lemma}
\newtheorem{proposition}[theorem]{Proposition}
\newtheorem{corollary}[theorem]{Corollary}

\theoremstyle{definition}
\newtheorem{definition}[theorem]{Definition}
\newtheorem{example}[theorem]{Example}
\newtheorem{construction}[theorem]{Construction}

\theoremstyle{remark}
\newtheorem{remark}[theorem]{Remark}
\newtheorem{notation}[theorem]{Notation}

% ── Python listings style ──────────────────────────────────────────────
\definecolor{jugeoBlue}{HTML}{2563EB}
\definecolor{jugeoGreen}{HTML}{059669}
\definecolor{jugeoRed}{HTML}{DC2626}
\definecolor{jugeoPurple}{HTML}{7C3AED}
\definecolor{jugeoGray}{HTML}{6B7280}
\definecolor{jugeoBg}{HTML}{F8FAFC}

\lstdefinestyle{jugeo-python}{
  language=Python,
  basicstyle=\ttfamily\small,
  keywordstyle=\color{jugeoBlue}\bfseries,
  stringstyle=\color{jugeoGreen},
  commentstyle=\color{jugeoGray}\itshape,
  numberstyle=\tiny\color{jugeoGray},
  backgroundcolor=\color{jugeoBg},
  numbers=left,
  numbersep=6pt,
  frame=single,
  framerule=0.4pt,
  rulecolor=\color{jugeoGray!40},
  breaklines=true,
  breakatwhitespace=true,
  showstringspaces=false,
  tabsize=4,
  xleftmargin=1.5em,
  framexleftmargin=1.5em,
  morekeywords={dataclass,frozen,field,Enum,IntEnum,auto,Any,
                Mapping,Sequence,Optional,match,case},
  literate={→}{{$\to$}}1 {←}{{$\leftarrow$}}1
           {≤}{{$\leq$}}1 {≥}{{$\geq$}}1 {≠}{{$\neq$}}1
           {∈}{{$\in$}}1 {∀}{{$\forall$}}1 {∃}{{$\exists$}}1
           {⊕}{{$\oplus$}}1 {⊖}{{$\ominus$}}1,
  escapeinside={(*@}{@*)},
}
\lstset{style=jugeo-python}

\lstdefinestyle{jugeo-output}{
  basicstyle=\ttfamily\small,
  backgroundcolor=\color{jugeoBg},
  frame=single,
  framerule=0.4pt,
  rulecolor=\color{jugeoGray!40},
  breaklines=true,
  showstringspaces=false,
  xleftmargin=1.5em,
  framexleftmargin=0.5em,
  numbers=none,
}

% ── JuGeo-specific macros ──────────────────────────────────────────────

% Judgment 8-tuple
\newcommand{\J}{\mathcal{J}}
\newcommand{\Jud}[8]{(#1,\,#2,\,#3,\,#4,\,#5,\,#6,\,#7,\,#8)}
\newcommand{\JudShort}[1]{\J_{#1}}

% Components
\newcommand{\coord}{\mathsf{c}}            % coordinate
\newcommand{\prop}{\varphi}                 % proposition
\newcommand{\carrier}{\mathcal{A}}          % carrier type
\newcommand{\evid}{\mathcal{E}}             % evidence bundle
\newcommand{\oblig}{\mathcal{O}}            % obligations
\newcommand{\obstr}{\mathcal{B}}            % obstructions
\newcommand{\trust}{\mathcal{T}}            % trust annotation
\newcommand{\prov}{\Pi}                     % provenance

% Trust algebra
\newcommand{\Talg}{\mathfrak{T}}            % trust ordered algebra
\newcommand{\Eadm}{\mathcal{E}_{\mathrm{adm}}}
\newcommand{\tleq}{\preceq}                 % trust ordering
\newcommand{\tjoin}{\oplus}                 % conservative join
\newcommand{\tatten}{\ominus}               % attenuation
\newcommand{\tpromote}{\uparrow_\pi}        % promotion
\newcommand{\tdemote}{\downarrow_\chi}       % demotion

% Trust tiers
\newcommand{\Tcontra}{\ensuremath{\mathsf{CONTRADICTED}}}
\newcommand{\Tunver}{\ensuremath{\mathsf{UNVERIFIED}}}
\newcommand{\Tcopilot}{\ensuremath{\mathsf{COPILOT}}}
\newcommand{\Toracle}{\ensuremath{\mathsf{ORACLE}}}
\newcommand{\Truntime}{\ensuremath{\mathsf{RUNTIME}}}
\newcommand{\Tsolver}{\ensuremath{\mathsf{SOLVER}}}
\newcommand{\Tproof}{\ensuremath{\mathsf{PROOF}}}

% Site and geometry
\newcommand{\Site}{\mathcal{S}}             % semantic site
\newcommand{\Cov}{\mathrm{Cov}}            % covering family
\newcommand{\Ob}{\mathrm{Ob}}              % objects
\newcommand{\Mor}{\mathrm{Mor}}            % morphisms
\newcommand{\res}{\mathrm{res}}            % restriction
\newcommand{\Cech}{\check{C}}              % Čech complex
\newcommand{\Hcoh}[1]{H^{#1}}             % cohomology group

% Descent
\newcommand{\descent}{\mathrm{desc}}
\newcommand{\glue}{\mathrm{glue}}
\newcommand{\overlap}[2]{#1 \cap #2}

% Semantic moves
\newcommand{\VERIFY}{\textsc{Verify}}
\newcommand{\CONSTRUCT}{\textsc{Construct}}
\newcommand{\NORMALIZE}{\textsc{Normalize}}
\newcommand{\GLUE}{\textsc{Glue}}
\newcommand{\RESTRICT}{\textsc{Restrict}}
\newcommand{\TRANSPORT}{\textsc{Transport}}
\newcommand{\REPAIR}{\textsc{Repair}}
\newcommand{\PROMOTE}{\textsc{Promote}}
\newcommand{\CHALLENGE}{\textsc{Challenge}}
\newcommand{\ESCALATE}{\textsc{Escalate}}

% Fragments
\newcommand{\QFLIA}{\texttt{QF\_LIA}}
\newcommand{\QFLRA}{\texttt{QF\_LRA}}
\newcommand{\QFBV}{\texttt{QF\_BV}}
\newcommand{\QFUF}{\texttt{QF\_UF}}

% Misc
\newcommand{\Python}{\textsc{Python}}
\newcommand{\JuGeo}{\textsc{JuGeo}}
\newcommand{\LEAN}{\textsc{Lean}\,4}
\newcommand{\Fstar}{F$^{\star}$}
\newcommand{\Coq}{\textsc{Coq}}
\newcommand{\Dafny}{\textsc{Dafny}}

% Common shortcuts
\newcommand{\ie}{i.e.\@\xspace}
\newcommand{\eg}{e.g.\@\xspace}
\newcommand{\cf}{cf.\@\xspace}
\newcommand{\wrt}{w.r.t.\@\xspace}

% Evaluation shortcuts
\newcommand{\numcases}{300}
\newcommand{\numspec}{100}
\newcommand{\numeq}{100}
\newcommand{\numbug}{100}
\newcommand{\medtime}{6.5\,ms}
\newcommand{\totaltime}{1.949\,s}


% ── Packages ────────────────────────────────────────────────────────────
\usepackage{amsmath,amssymb,amsthm,mathtools}
\usepackage{tikz-cd}
\usepackage{listings}
\usepackage{xcolor}
\usepackage{xspace}
\usepackage{booktabs}
\usepackage{multirow}
\usepackage{enumitem}
\usepackage[expansion=false]{microtype}
\usepackage{hyperref}
\usepackage{cleveref}
% \usepackage{thmtools}  % disabled: not available in this TeX installation
\usepackage{float}
\usepackage{caption}
\usepackage{subcaption}
\usepackage{tabularx}
\usepackage{stmaryrd}       % \llbracket, \rrbracket
\usepackage{bbm}            % \mathbbm{1}

% ── Page geometry ───────────────────────────────────────────────────────
\usepackage[margin=1in]{geometry}

% ── Theorem environments ────────────────────────────────────────────────
\theoremstyle{plain}
\newtheorem{theorem}{Theorem}[section]
\newtheorem{lemma}[theorem]{Lemma}
\newtheorem{proposition}[theorem]{Proposition}
\newtheorem{corollary}[theorem]{Corollary}

\theoremstyle{definition}
\newtheorem{definition}[theorem]{Definition}
\newtheorem{example}[theorem]{Example}
\newtheorem{construction}[theorem]{Construction}

\theoremstyle{remark}
\newtheorem{remark}[theorem]{Remark}
\newtheorem{notation}[theorem]{Notation}

% ── Python listings style ──────────────────────────────────────────────
\definecolor{jugeoBlue}{HTML}{2563EB}
\definecolor{jugeoGreen}{HTML}{059669}
\definecolor{jugeoRed}{HTML}{DC2626}
\definecolor{jugeoPurple}{HTML}{7C3AED}
\definecolor{jugeoGray}{HTML}{6B7280}
\definecolor{jugeoBg}{HTML}{F8FAFC}

\lstdefinestyle{jugeo-python}{
  language=Python,
  basicstyle=\ttfamily\small,
  keywordstyle=\color{jugeoBlue}\bfseries,
  stringstyle=\color{jugeoGreen},
  commentstyle=\color{jugeoGray}\itshape,
  numberstyle=\tiny\color{jugeoGray},
  backgroundcolor=\color{jugeoBg},
  numbers=left,
  numbersep=6pt,
  frame=single,
  framerule=0.4pt,
  rulecolor=\color{jugeoGray!40},
  breaklines=true,
  breakatwhitespace=true,
  showstringspaces=false,
  tabsize=4,
  xleftmargin=1.5em,
  framexleftmargin=1.5em,
  morekeywords={dataclass,frozen,field,Enum,IntEnum,auto,Any,
                Mapping,Sequence,Optional,match,case},
  literate={→}{{$\to$}}1 {←}{{$\leftarrow$}}1
           {≤}{{$\leq$}}1 {≥}{{$\geq$}}1 {≠}{{$\neq$}}1
           {∈}{{$\in$}}1 {∀}{{$\forall$}}1 {∃}{{$\exists$}}1
           {⊕}{{$\oplus$}}1 {⊖}{{$\ominus$}}1,
  escapeinside={(*@}{@*)},
}
\lstset{style=jugeo-python}

\lstdefinestyle{jugeo-output}{
  basicstyle=\ttfamily\small,
  backgroundcolor=\color{jugeoBg},
  frame=single,
  framerule=0.4pt,
  rulecolor=\color{jugeoGray!40},
  breaklines=true,
  showstringspaces=false,
  xleftmargin=1.5em,
  framexleftmargin=0.5em,
  numbers=none,
}

% ── JuGeo-specific macros ──────────────────────────────────────────────

% Judgment 8-tuple
\newcommand{\J}{\mathcal{J}}
\newcommand{\Jud}[8]{(#1,\,#2,\,#3,\,#4,\,#5,\,#6,\,#7,\,#8)}
\newcommand{\JudShort}[1]{\J_{#1}}

% Components
\newcommand{\coord}{\mathsf{c}}            % coordinate
\newcommand{\prop}{\varphi}                 % proposition
\newcommand{\carrier}{\mathcal{A}}          % carrier type
\newcommand{\evid}{\mathcal{E}}             % evidence bundle
\newcommand{\oblig}{\mathcal{O}}            % obligations
\newcommand{\obstr}{\mathcal{B}}            % obstructions
\newcommand{\trust}{\mathcal{T}}            % trust annotation
\newcommand{\prov}{\Pi}                     % provenance

% Trust algebra
\newcommand{\Talg}{\mathfrak{T}}            % trust ordered algebra
\newcommand{\Eadm}{\mathcal{E}_{\mathrm{adm}}}
\newcommand{\tleq}{\preceq}                 % trust ordering
\newcommand{\tjoin}{\oplus}                 % conservative join
\newcommand{\tatten}{\ominus}               % attenuation
\newcommand{\tpromote}{\uparrow_\pi}        % promotion
\newcommand{\tdemote}{\downarrow_\chi}       % demotion

% Trust tiers
\newcommand{\Tcontra}{\ensuremath{\mathsf{CONTRADICTED}}}
\newcommand{\Tunver}{\ensuremath{\mathsf{UNVERIFIED}}}
\newcommand{\Tcopilot}{\ensuremath{\mathsf{COPILOT}}}
\newcommand{\Toracle}{\ensuremath{\mathsf{ORACLE}}}
\newcommand{\Truntime}{\ensuremath{\mathsf{RUNTIME}}}
\newcommand{\Tsolver}{\ensuremath{\mathsf{SOLVER}}}
\newcommand{\Tproof}{\ensuremath{\mathsf{PROOF}}}

% Site and geometry
\newcommand{\Site}{\mathcal{S}}             % semantic site
\newcommand{\Cov}{\mathrm{Cov}}            % covering family
\newcommand{\Ob}{\mathrm{Ob}}              % objects
\newcommand{\Mor}{\mathrm{Mor}}            % morphisms
\newcommand{\res}{\mathrm{res}}            % restriction
\newcommand{\Cech}{\check{C}}              % Čech complex
\newcommand{\Hcoh}[1]{H^{#1}}             % cohomology group

% Descent
\newcommand{\descent}{\mathrm{desc}}
\newcommand{\glue}{\mathrm{glue}}
\newcommand{\overlap}[2]{#1 \cap #2}

% Semantic moves
\newcommand{\VERIFY}{\textsc{Verify}}
\newcommand{\CONSTRUCT}{\textsc{Construct}}
\newcommand{\NORMALIZE}{\textsc{Normalize}}
\newcommand{\GLUE}{\textsc{Glue}}
\newcommand{\RESTRICT}{\textsc{Restrict}}
\newcommand{\TRANSPORT}{\textsc{Transport}}
\newcommand{\REPAIR}{\textsc{Repair}}
\newcommand{\PROMOTE}{\textsc{Promote}}
\newcommand{\CHALLENGE}{\textsc{Challenge}}
\newcommand{\ESCALATE}{\textsc{Escalate}}

% Fragments
\newcommand{\QFLIA}{\texttt{QF\_LIA}}
\newcommand{\QFLRA}{\texttt{QF\_LRA}}
\newcommand{\QFBV}{\texttt{QF\_BV}}
\newcommand{\QFUF}{\texttt{QF\_UF}}

% Misc
\newcommand{\Python}{\textsc{Python}}
\newcommand{\JuGeo}{\textsc{JuGeo}}
\newcommand{\LEAN}{\textsc{Lean}\,4}
\newcommand{\Fstar}{F$^{\star}$}
\newcommand{\Coq}{\textsc{Coq}}
\newcommand{\Dafny}{\textsc{Dafny}}

% Common shortcuts
\newcommand{\ie}{i.e.\@\xspace}
\newcommand{\eg}{e.g.\@\xspace}
\newcommand{\cf}{cf.\@\xspace}
\newcommand{\wrt}{w.r.t.\@\xspace}

% Evaluation shortcuts
\newcommand{\numcases}{300}
\newcommand{\numspec}{100}
\newcommand{\numeq}{100}
\newcommand{\numbug}{100}
\newcommand{\medtime}{6.5\,ms}
\newcommand{\totaltime}{1.949\,s}

% \input{data-paper89}  % data file removed: paper 89 has no measured data
\usepackage{array}
\usepackage{tikz}
\usetikzlibrary{arrows.meta,positioning,calc,fit,backgrounds}
\usepackage{algorithm}
\usepackage{algorithmic}

% ── paper-specific macros ─────────────────────────────────────────
\newcommand{\CopChan}{\texttt{CopilotOracleChannel}}
\newcommand{\CopOracle}{\texttt{CopilotOracle}}
\newcommand{\CopFallback}{\texttt{CopilotFallbackPolicy}}
\newcommand{\CopFragment}{\texttt{CopilotFragmentAssist}}
\newcommand{\OChan}{\texttt{OracleChannel}}
\newcommand{\OJuris}{\texttt{OracleJurisdiction}}
\newcommand{\TCeil}{\texttt{TrustCeilingEnforcer}}
\newcommand{\OPRec}{\texttt{OracleProposalRecord}}
\newcommand{\copilottrust}{$\mathsf{copilot\_suggested}$}
\newcommand{\ceilop}{\mathsf{ceil}}
\newcommand{\promote}{\mathsf{promote}}
\newcommand{\corroborate}{\mathsf{corroborate}}
\newcommand{\selfpromo}{\mathsf{self\text{-}promo}}

\title{\textbf{Copilot as Federated Oracle:\\
LLMs in the Solver Network}\\[6pt]
{\large Paper~89 of the Judgment Geometry Series}}

\author{Halley Young \\
\textit{Microsoft Research} \\
\texttt{halley.young@microsoft.com}}

\date{\today}

\begin{document}
\maketitle

\noindent\textbf{Audit note.}
The repository contains oracle-federation infrastructure such as
\texttt{OracleProposalRecord}, \texttt{OracleJurisdiction}, and the
controlled-oracle machinery, but not every copilot-specific class shown
in the original draft is exported as a stable JuGeo API.  In this
audited version, any unmatched listings should be read as simplified
interface sketches around the real trust-ceiling and jurisdiction logic.

% ══════════════════════════════════════════════════════════════════
%  Abstract
% ══════════════════════════════════════════════════════════════════

\begin{abstract}
Modern software verification systems increasingly rely on heterogeneous
solver backends---SMT solvers, model checkers, proof assistants, and
now large language models (LLMs) such as GitHub Copilot.  Integrating
an LLM into a \emph{trusted} verification pipeline poses unique
challenges: the model can fabricate evidence, self-promote its own
confidence, and drift outside its jurisdiction.  We present a
\emph{federated oracle} architecture within the Judgment Geometry
framework that treats Copilot as a \emph{controlled oracle channel}
with a hard trust ceiling at \copilottrust{}.  Our design enforces
three invariants: (1)~\emph{no self-promotion}---the oracle cannot
raise its own trust tier; (2)~\emph{fallback soundness}---copilot is
consulted only when primary backends are exhausted; and
(3)~\emph{corroboration necessity}---any suggestion promoted above
the ceiling must be independently corroborated.  The repository-backed
core of this story is the controlled-oracle and trust-ceiling
infrastructure; copilot-specific policy layers shown below are best read
as design sketches, and earlier claims about Lean mechanization and
measured multi-domain evaluation have been removed.
\end{abstract}

% ══════════════════════════════════════════════════════════════════
%  1. Introduction
% ══════════════════════════════════════════════════════════════════

\section{Introduction}\label{sec:intro}

The Judgment Geometry framework models software artefacts through
\emph{judgments}---triples $\Jud{J} = (\coord, \prop, \evid)$ whose
coordinates locate a code entity, whose propositions describe its
intended behaviour, and whose evidence witnesses the verification
status of that behaviour~\cite{young2024jg1,young2024jg4}.  Each
judgment carries a \emph{trust level}
$\tau \in \{\Tcontra, \Tunver, \Tcopilot, \Toracle, \Truntime,
\Tsolver, \Tproof\}$
ordered by the lattice $(\mathcal{T}, \tleq)$.

Over the preceding 88~papers in this series we have developed the
algebraic, topological, and categorical infrastructure for composing
verification evidence across heterogeneous backends.  Papers~7
and~8~\cite{young2024jg7,young2024jg8} introduced the solver
router---a central dispatcher that directs verification requests to
the appropriate backend.  Paper~31~\cite{young2024jg31} formalised
oracle channels as trust-bounded interfaces for querying external
knowledge sources.

In this paper we ask: \emph{What happens when one of those oracles
is an LLM?}  Large language models such as GitHub Copilot have become
indispensable coding assistants, yet they can hallucinate proofs,
fabricate counterexamples, and---most dangerously---present
high-confidence claims that escape the trust discipline of the
verification pipeline.  Naively plugging an LLM into the solver
network would compromise the soundness guarantees that Judgment
Geometry is designed to uphold.

Our solution is a \emph{federated oracle} architecture that treats
Copilot as a first-class but \emph{trust-bounded} participant in the
solver network.  The key contributions are:

\begin{enumerate}
  \item \textbf{Controlled copilot-oracle architecture.}  A trust-bounded
    channel design centered on repository-backed oracle-jurisdiction and
    proposal-record ideas, with copilot-specific methods presented as
    interface sketches (Section~\ref{sec:channel}).

  \item \textbf{Fallback and fragment-assist policies.}
    Policy sketches governing when a copilot oracle would be consulted
    and how it might assist fragment classification
    (Section~\ref{sec:fallback}).

  \item \textbf{Formal properties.}  Three desired invariants---no
    self-promotion, fallback soundness, and corroboration
    necessity---stated mathematically rather than claimed as shipped
    Lean artifacts (Section~\ref{sec:formal}).

  \item \textbf{Evaluation methodology.}  An experimental
    framework for measuring suggestion acceptance, corroboration
    effectiveness, and self-promotion violation detection
    (Section~\ref{sec:experiments}).
\end{enumerate}

The challenge is not merely engineering.  An LLM oracle introduces
a qualitatively new failure mode: \emph{trust inflation}.  Unlike
an SMT solver, which either returns $\mathsf{sat}$/$\mathsf{unsat}$
or times out, an LLM can emit natural-language prose that
\emph{appears} to carry higher epistemic authority than it warrants.
A response like ``this lemma is trivially true by induction'' may
be syntactically well-formed and even correct, yet assigning it
$\Tsolver$-level trust would be unsound.  Our architecture
addresses this by enforcing a hard ceiling at $\Tcopilot$ and
requiring independent corroboration for any promotion.

The design philosophy follows the \emph{principle of least
authority} (POLA): the copilot oracle receives exactly the
permissions it needs to be useful---the ability to suggest
encodings, explanations, and decompositions---and no more.  It
cannot modify the trust lattice, cannot promote its own output,
and cannot bypass the jurisdiction check.

The rest of this paper is organised as follows.
Section~\ref{sec:federation} reviews the oracle federation model in
Judgment Geometry.  Section~\ref{sec:channel} presents the
\CopChan{} class.  Section~\ref{sec:fallback} describes fallback
and fragment-assist policies.  Section~\ref{sec:formal} states and
proves the three core theorems.  Section~\ref{sec:experiments}
reports experimental results.  Section~\ref{sec:related} surveys
related work, and Section~\ref{sec:conclusion} concludes.

% ══════════════════════════════════════════════════════════════════
%  2. Oracle Federation in Judgment Geometry
% ══════════════════════════════════════════════════════════════════

\section{Oracle Federation in Judgment Geometry}
\label{sec:federation}

The oracle federation model in Judgment Geometry provides a
principled framework for integrating external knowledge sources
into the verification pipeline.  Each oracle is wrapped in a
\emph{channel} that enforces jurisdiction, trust ceilings, and
audit logging.

The federation model is motivated by the observation that modern
verification systems rarely rely on a single backend.  A typical
verification task might require an SMT solver for arithmetic
constraints, a model checker for temporal properties, a proof
assistant for complex inductive arguments, and now an LLM for
heuristic guidance.  Each of these backends has different strengths,
different failure modes, and different levels of trustworthiness.

The oracle federation provides a uniform interface---the
$\mathsf{OracleChannel}$---that abstracts over these differences
while preserving the trust discipline.  Every channel declares its
jurisdiction (what kinds of queries it can handle), its trust
ceiling (the maximum trust level its responses can carry), and its
audit policy (how interactions are logged for post-hoc review).

\subsection{Controlled Oracles}\label{sec:controlled}

An oracle in Judgment Geometry is any external agent that can
supply evidence for a judgment.  The $\mathsf{OracleChannel}$ base
class (from \texttt{jugeo.foundations.oracle\_federation.controlled%
\_oracles}) manages this interaction through a structured protocol:

\begin{enumerate}
  \item The solver router selects a channel based on the query
    domain.
  \item The channel checks jurisdiction (\ie{} whether the query
    falls within the oracle's declared competence).
  \item If accepted, the channel invokes the oracle and wraps the
    response in an \OPRec{}.
  \item A $\mathsf{TrustCeilingEnforcer}$ clamps the trust level
    of the response to the channel's ceiling.
\end{enumerate}

The key invariant is that no oracle can unilaterally raise its own
trust level beyond its assigned ceiling.  Formally, for any oracle
channel $c$ with ceiling $\kappa_c$ and any response $r$ produced
by $c$, we require:
\begin{equation}\label{eq:ceiling}
  \tau(r) \;\tleq\; \kappa_c.
\end{equation}

This invariant is the cornerstone of the entire oracle federation.
Without it, a single compromised or misconfigured oracle could
propagate inflated trust levels through the verification pipeline,
undermining the guarantees of every downstream consumer.  The
invariant is enforced at two levels: locally by each channel's
$\mathsf{propose}$ method, and globally by the
$\mathsf{TrustCeilingEnforcer}$ that sits in the pipeline between
the channel and the proposal log.

\begin{definition}[Oracle Channel]\label{def:oracle-channel}
An \emph{oracle channel} $c = (id, \mathcal{J}, \kappa, f)$
consists of a unique identifier $id$, a jurisdiction $\mathcal{J}$,
a trust ceiling $\kappa \in \mathcal{T}$, and an invocation function
$f : \mathsf{Request} \to \mathsf{Response}$.  The channel's
$\mathsf{propose}$ method computes $f(q)$ for a request $q$ within
$\mathcal{J}$ and clamps the result's trust to $\kappa$.
\end{definition}

\begin{lstlisting}[style=jugeo-python,caption={Simplified channel sketch.}]
# Schematic interface sketch.  The repository contains controlled-oracle
# machinery, but not every symbol in this listing is exported exactly as
# shown.

class OracleChannel:
    """Trust-bounded oracle channel (Theory2 S7.1.1)."""

    def __init__(self, oracle_id=None, name="oracle",
                 channel_kind="oracle",
                 jurisdiction=None,
                 trust_ceiling="oracle_proposed"):
        self._oracle_id = oracle_id or uuid.uuid4().hex[:16]
        self._name = name
        self._channel_kind = channel_kind
        self._jurisdiction = jurisdiction
        self._trust_ceiling = trust_ceiling
        self._proposal_log: list[OracleProposalRecord] = []

    def propose(self, request):
        """Submit request; return clamped response dict."""
        if self._jurisdiction and not \
                self._jurisdiction.check_request(request):
            raise JurisdictionError(request)
        raw = self._invoke_oracle(request)
        return self._clamp(raw)
\end{lstlisting}

The $\mathsf{propose}$ method is the sole entry point for oracle
interaction.  It delegates to a private $\mathsf{\_invoke\_oracle}$
method (overridden by subclasses) and then applies the trust ceiling
via $\mathsf{\_clamp}$.

\subsection{Oracle Jurisdiction}\label{sec:jurisdiction}

Each oracle channel declares an \OJuris{} that specifies its
authorised domains.  This prevents, for example, a Copilot channel
from answering queries about bitvector overflow when it has only
been authorised for string-theory suggestions.

\begin{definition}[Oracle Jurisdiction]\label{def:jurisdiction}
An \emph{oracle jurisdiction} $\mathcal{J} = (s, D, \kappa)$
consists of a scope identifier $s$, a set of allowed domains
$D \subseteq \mathcal{D}$, and a trust ceiling $\kappa$.  A
request $q$ is \emph{within jurisdiction} if
$\mathrm{domain}(q) \in D$.
\end{definition}

The jurisdiction check is enforced before any oracle invocation.
If a request falls outside jurisdiction, the channel raises a
$\mathsf{JurisdictionError}$ and no proposal record is created.

\begin{lstlisting}[style=jugeo-python,caption={OracleJurisdiction class.}]
class OracleJurisdiction:
    """Declares authorized domains for oracle."""

    def __init__(self, scope, allowed_domains,
                 trust_ceiling="oracle_proposed",
                 description=""):
        self.scope = scope
        self.allowed_domains = set(allowed_domains)
        self.trust_ceiling = trust_ceiling
        self.description = description

    def check_request(self, request_kind):
        """Return True iff request_kind is in scope."""
        return request_kind in self.allowed_domains
\end{lstlisting}

A crucial design property is that jurisdictions of distinct oracle
channels should be \emph{disjoint} or \emph{hierarchically nested}:

\begin{proposition}[Jurisdiction Disjointness]\label{prop:disjoint}
If oracle channels $c_1$ and $c_2$ have jurisdictions
$\mathcal{J}_1 = (s_1, D_1, \kappa_1)$ and
$\mathcal{J}_2 = (s_2, D_2, \kappa_2)$ with $s_1 \ne s_2$,
then their responses never conflict on the same query domain:
for any request $q$ within both jurisdictions, the router
selects the channel with higher ceiling.
\end{proposition}

This property prevents ambiguity in the routing layer.  When the
solver router receives a verification request, it consults the
jurisdiction registry to determine which channels are eligible.  If
two channels claim authority over the same domain, the router must
resolve the conflict.  Disjointness eliminates this case entirely;
hierarchical nesting resolves it by preferring the more specific
jurisdiction.

In the copilot setting, the copilot channel's jurisdiction is
typically a subset of the broader oracle jurisdiction, restricted
to domains where LLM heuristics are known to be helpful (\eg{}
$\QFLIA$ encoding suggestions, string theory explanations).  The
copilot channel never has jurisdiction over domains where formal
solvers provide definitive answers (\eg{} $\QFBV$ satisfiability).

\subsection{Trust Ceiling Enforcement}\label{sec:enforcement}

The \TCeil{} is a policy guard that sits between the oracle
channel and the proposal log.  It maintains a registry of
channel-to-ceiling mappings and clamps every response.

\begin{lstlisting}[style=jugeo-python,caption={TrustCeilingEnforcer.}]
class TrustCeilingEnforcer:
    """Policy guard clamping oracle responses."""

    def __init__(self):
        self.ceiling_map: dict[str, str] = {}
        self.violation_log: list[dict] = []

    def register_ceiling(self, channel_id, ceiling_level):
        self.ceiling_map[channel_id] = ceiling_level

    def enforce(self, channel_id, response_dict):
        ceiling = self.ceiling_map.get(channel_id,
                                       "oracle_proposed")
        actual = response_dict.get("trust_level",
                                    "unverified")
        if _trust_ord(actual) > _trust_ord(ceiling):
            self.violation_log.append({
                "channel": channel_id,
                "requested": actual,
                "clamped_to": ceiling,
            })
            response_dict["trust_level"] = ceiling
        return response_dict
\end{lstlisting}

The enforcer operates as a pure function on response dictionaries:
it never modifies the oracle itself, only the outgoing trust tag.
This separation of concerns ensures that enforcement is
\emph{compositional}---multiple enforcers can be stacked in a
pipeline without interference.

The compositionality of ceiling enforcement is essential for the
federated setting.  When multiple oracle channels are active
simultaneously, each with its own ceiling, the global enforcer
applies ceilings independently.  The resulting trust assignment is
the pointwise minimum of the channel-specific ceiling and the
claimed trust level.

Formally, let $E$ be the enforcer with ceiling map
$\{c_i \mapsto \kappa_i\}_{i=1}^{n}$.  For a response $r$ from
channel $c_j$ with claimed trust $\tau(r)$, the enforced trust is:
\begin{equation}\label{eq:enforced}
  E(c_j, r) \;=\; \min\bigl(\tau(r),\; \kappa_j\bigr).
\end{equation}

This is monotone: raising $\kappa_j$ can only raise (never lower)
the enforced trust, and lowering $\kappa_j$ can only lower it.  The
enforcer thus acts as a \emph{Galois connection} between the lattice
of claimed trust levels and the lattice of enforced trust levels.

\begin{remark}
The violation log maintained by $\mathsf{TrustCeilingEnforcer}$ is
append-only and serves as an audit trail.  In production, this log
is forwarded to the system's provenance store for post-hoc analysis.
\end{remark}

% ══════════════════════════════════════════════════════════════════
%  3. The CopilotOracleChannel
% ══════════════════════════════════════════════════════════════════

\section{The CopilotOracleChannel}\label{sec:channel}

The \CopChan{} class extends \OChan{} with Copilot-specific
semantics.  Its trust ceiling is permanently fixed at
\copilottrust{}, reflecting the epistemic status of LLM-generated
content: useful but unverified.

Figure~\ref{fig:channel-arch} illustrates the architecture.  A
query enters through the solver router, is dispatched to the
\CopChan{}, and passes through jurisdiction checking, the
suggestion engine, self-promotion validation, and ceiling
enforcement before a proposal record is emitted.

\begin{figure}[ht]
\centering
\begin{tikzpicture}[
  node distance=1.2cm and 1.8cm,
  box/.style={draw, rounded corners, minimum width=2.8cm,
              minimum height=0.8cm, align=center,
              font=\small},
  arr/.style={-{Stealth[length=3mm]}, thick},
]
  \node[box, fill=blue!10] (router) {Solver\\Router};
  \node[box, fill=green!10, right=of router] (juris)
    {Jurisdiction\\Check};
  \node[box, fill=orange!10, right=of juris] (suggest)
    {Copilot\\Suggest};
  \node[box, fill=red!10, below=of suggest] (selfpromo)
    {Self-Promotion\\Validator};
  \node[box, fill=purple!10, left=of selfpromo] (enforcer)
    {Ceiling\\Enforcer};
  \node[box, fill=gray!10, left=of enforcer] (log)
    {Proposal\\Log};

  \draw[arr] (router) -- (juris);
  \draw[arr] (juris) -- (suggest);
  \draw[arr] (suggest) -- (selfpromo);
  \draw[arr] (selfpromo) -- (enforcer);
  \draw[arr] (enforcer) -- (log);

  \draw[arr, dashed, red] (juris.south) --
    ++(0,-0.5) node[right, font=\scriptsize, red]
    {refuse};
  \draw[arr, dashed, red] (selfpromo.west) --
    ++(-0.4,0) node[above, font=\scriptsize, red]
    {violation};
\end{tikzpicture}
\caption{Architecture of the \CopChan{}.  Dashed arrows indicate
  refusal paths.}
\label{fig:channel-arch}
\end{figure}

\subsection{The \texorpdfstring{$\mathsf{suggest}$}{suggest} Method}
\label{sec:suggest}

The $\mathsf{suggest}$ method is the primary interface through which
the solver router obtains Copilot's input.  It accepts a context
dictionary describing the current verification goal and returns a
suggestion record.

\begin{lstlisting}[style=jugeo-python,caption={\CopChan{} sketch (core interface).}]
class CopilotOracleChannel(OracleChannel):
    """Copilot-specific oracle channel."""

    def __init__(self, oracle_id=None,
                 name="copilot",
                 trust_ceiling="copilot_suggested",
                 jurisdiction=None):
        super().__init__(
            oracle_id=oracle_id,
            name=name,
            channel_kind="copilot",
            jurisdiction=jurisdiction,
            trust_ceiling=trust_ceiling,
        )
        self.suggestion_history: list[dict] = []
        self.refusal_reasons: list[dict] = []

    def suggest(self, context):
        """Generate suggestion for verification context."""
        self._validate_jurisdiction(context)
        suggestion = self._invoke_copilot(context)
        self.suggestion_history.append({
            "context": context,
            "suggestion": suggestion,
            "trust": self._trust_ceiling,
        })
        return suggestion

    def suggestion_count(self):
        """Return total suggestions emitted."""
        return len(self.suggestion_history)
\end{lstlisting}

The $\mathsf{suggest}$ method never returns a trust level above
\copilottrust{}.  This is enforced both by the method itself (which
tags every suggestion at the channel's ceiling) and by the external
$\mathsf{TrustCeilingEnforcer}$ (which clamps any residual
over-claims).

The dual enforcement---local clamping within the channel and global
clamping by the enforcer---provides defence in depth.  Even if a
subclass override inadvertently emits an inflated trust tag, the
external enforcer catches it.  This layered design is motivated by
the observation that LLM outputs are inherently unpredictable: a
model might embed trust-level metadata in its response that bypasses
the channel's local clamping logic.

Each entry in $\mathsf{suggestion\_history}$ records the context,
the suggestion, and the trust level at emission.  This log enables
post-hoc auditing: one can verify that every suggestion was indeed
capped at the ceiling by replaying the log through the enforcer.

The $\mathsf{suggestion\_count}$ method provides a lightweight
query over the history for rate-limiting and monitoring purposes.

\subsection{The \texorpdfstring{$\mathsf{refuse}$}{refuse} Method}
\label{sec:refuse}

When a query falls outside the channel's jurisdiction or when the
self-promotion validator detects a violation, the channel refuses
to produce a suggestion.  The $\mathsf{refuse}$ method logs the
reason and returns a null response.

\begin{lstlisting}[style=jugeo-python,caption={Refusal logic.}]
    def refuse(self, reason, context):
        """Refuse to answer; log reason."""
        self.refusal_reasons.append({
            "reason": reason,
            "context": context,
        })
        return None
\end{lstlisting}

Refusals are first-class events in the audit trail.  The
evaluation framework in Section~\ref{sec:experiments} measures
refusals broken down by cause: jurisdiction violations, ceiling
violations, and self-promotion attempts.

\subsection{The \texorpdfstring{$\mathsf{corroborate}$}{corroborate}
  Method}\label{sec:corroborate}

A copilot suggestion that the solver wants to promote above
\copilottrust{} must be independently corroborated.  The
$\mathsf{corroborate}$ method records such corroboration:

\begin{lstlisting}[style=jugeo-python,caption={Corroboration protocol.}]
    def corroborate(self, proposal_id, evidence):
        """Record external corroboration for a proposal."""
        rec = self._find_proposal(proposal_id)
        if rec is None:
            raise ValueError(f"Unknown proposal: {proposal_id}")
        rec.corroborate(evidence["source_id"],
                        evidence["trust_level"])
        return rec
\end{lstlisting}

The corroboration source must satisfy two conditions:
\begin{enumerate}
  \item \emph{Independence}: the source must not be the copilot
    channel itself (enforced by checking
    $\mathsf{source\_id} \ne \mathsf{oracle\_id}$).
  \item \emph{Sufficient trust}: the source's trust level must
    meet or exceed the target promotion level.
\end{enumerate}

The corroboration success rate is measured by the evaluation
framework in Section~\ref{sec:experiments}.

The corroboration protocol interacts with the \OPRec{} class, which
maintains a list of corroborating sources for each proposal.  When
a corroboration is recorded, the proposal record is updated with
the source's identifier and trust level.  This allows post-hoc
analysis of which corroboration sources were most effective for
different query domains.

\begin{remark}
The independence requirement for corroboration sources prevents a
subtle attack vector: if the copilot could corroborate its own
suggestions (even through an alias), it could bypass the trust
ceiling by first emitting a suggestion and then corroborating it.
The $\mathsf{source\_id} \ne \mathsf{oracle\_id}$ check in the
$\mathsf{corroborate}$ method prevents this.  This is formalised
in Lean~4 as $\mathsf{corroboration\_source\_independent}$.
\end{remark}

\subsection{The \texorpdfstring{$\mathsf{validate\_no\_self\_promotion}$}%
  {validate\_no\_self\_promotion} Method}\label{sec:selfpromo}

The most critical safety check is the self-promotion validator.
An LLM, left unchecked, could embed claims like ``this proof is
mechanically verified'' in its output, causing downstream components
to assign an inflated trust level.

\begin{definition}[Self-Promotion]\label{def:selfpromo}
A suggestion $s$ from oracle channel $c$ with ceiling $\kappa_c$
is \emph{self-promoting} if $s$ contains metadata or natural-language
assertions that, if taken at face value, would assign $s$ a trust
level $\tau > \kappa_c$.
\end{definition}

Self-promotion is subtle because it can take many forms.  A direct
metadata override (setting $\mathsf{trust\_level}$ to
$\mathsf{mechanically\_verified}$) is easy to detect.  More
insidious is \emph{implicit} self-promotion, where the suggestion's
natural-language content contains phrases that downstream
NLP-based components might parse as trust assertions.  For example,
a copilot response containing ``I have formally verified this
lemma'' could trigger an automated trust-extraction pipeline to
assign $\Tproof$-level trust.

The $\mathsf{validate\_no\_self\_promotion}$ method scans the
channel's suggestion history for self-promoting content.  It
operates in three phases:

\begin{enumerate}
  \item \textbf{Metadata scan.}  Check whether any suggestion's
    $\mathsf{trust\_level}$ field exceeds the channel ceiling.
  \item \textbf{Keyword scan.}  Search suggestion text for
    trust-inflating keywords (\eg{} ``verified'', ``proved'',
    ``mechanically checked'').
  \item \textbf{Structural scan.}  Detect embedded JSON or
    structured metadata that could be parsed as a higher trust
    assertion by downstream consumers.
\end{enumerate}

\begin{remark}
The keyword scan is necessarily heuristic.  However, because
the \TCeil{} provides a hard backstop, false negatives in the
keyword scan do not compromise soundness---they only reduce the
quality of the audit trail.
\end{remark}

% ══════════════════════════════════════════════════════════════════
%  4. Solver Fallback and Fragment Assist
% ══════════════════════════════════════════════════════════════════

\section{Solver Fallback and Fragment Assist}\label{sec:fallback}

The copilot oracle is not the first choice for any verification
query---it is a \emph{fallback} consulted only when primary
backends have been exhausted or when the query requires heuristic
guidance that formal solvers cannot provide.  This section describes
the two policy classes that govern copilot's role in the solver
network.

\subsection{CopilotFallbackPolicy}\label{sec:fallback-policy}

The \CopFallback{} class encapsulates the decision procedure for
when to consult the copilot oracle.  It is initialised with four
parameters:

\begin{itemize}
  \item $\mathsf{\_enabled}$: global kill switch.
  \item $\mathsf{\_max\_qpm}$: rate limit (queries per minute).
  \item $\mathsf{\_corroboration\_threshold}$: the trust tier
    above which corroboration is mandatory.
  \item $\mathsf{\_ceiling}$: the effective trust ceiling for
    copilot responses.
\end{itemize}

\begin{lstlisting}[style=jugeo-python,caption={Copilot fallback-policy sketch.}]
# Design sketch only: the exact fallback-policy API below is not a
# stable repository export.

class CopilotFallbackPolicy:
    """Policy: when to consult the copilot oracle."""

    def __init__(self, *, enabled=True,
                 max_queries_per_minute=10,
                 require_corroboration_above=
                     TrustTier.PROPOSAL,
                 copilot_trust_ceiling=
                     TrustTier.PROPOSAL):
        self._enabled = enabled
        self._max_qpm = max_queries_per_minute
        self._corroboration_threshold = \
            require_corroboration_above
        self._ceiling = copilot_trust_ceiling
        self._query_timestamps: list[float] = []

    def when_to_use_copilot(self, domains,
                            other_backends_exhausted):
        """Decide whether to fall back to copilot."""
        if not self._enabled:
            return False
        if not self.rate_limit():
            return False
        if other_backends_exhausted:
            return True
        return self._heuristic_match(domains)

    def trust_ceiling_for_request(self, requested_tier):
        """Return effective ceiling for request."""
        return min(requested_tier, self._ceiling)

    def require_corroboration(self, effective_tier):
        """Is corroboration needed at this tier?"""
        return effective_tier > self._corroboration_threshold

    def rate_limit(self):
        """Return True if under rate limit."""
        now = time.time()
        cutoff = now - 60.0
        self._query_timestamps = [
            t for t in self._query_timestamps
            if t > cutoff
        ]
        if len(self._query_timestamps) >= self._max_qpm:
            return False
        self._query_timestamps.append(now)
        return True
\end{lstlisting}

The $\mathsf{when\_to\_use\_copilot}$ method implements a
two-phase decision:
\begin{enumerate}
  \item \emph{Guard phase}: check enabled flag and rate limit.
  \item \emph{Dispatch phase}: if other backends are exhausted,
    always consult copilot; otherwise apply a heuristic domain
    match.
\end{enumerate}

\begin{table}[ht]
\centering
\caption{Fallback policy parameters and experimental values.}
\label{tab:fallback-params}
\begin{tabular}{lcc}
\hline
\textbf{Parameter} & \textbf{Type} & \textbf{Value} \\
\hline
$\mathsf{enabled}$ & $\mathsf{bool}$ & $\mathsf{True}$ \\
$\mathsf{max\_qpm}$ & $\mathsf{int}$ & 10 \\
$\mathsf{corroboration\_threshold}$ & $\mathsf{TrustTier}$
  & $\mathsf{PROPOSAL}$ \\
$\mathsf{ceiling}$ & $\mathsf{TrustTier}$
  & $\mathsf{PROPOSAL}$ \\
\hline
\end{tabular}
\end{table}

The fallback policy records all invocations and tracks the
fraction that occur when all primary backends have been exhausted
versus those triggered by heuristic domain matching.

\subsection{CopilotFragmentAssist}\label{sec:fragment-assist}

The \CopFragment{} class provides LLM-assisted guidance for the
fragment classifier.  SMT-LIB theory fragments
($\QFLIA$, $\QFLRA$, $\QFBV$, $\QFUF$, \etc) determine which
solver backend is appropriate for a given formula.  Fragment
classification is typically syntactic, but borderline cases---especially
mixed-theory formulas---benefit from heuristic assistance.

\CopFragment{} wraps a $\mathsf{FragmentClassifier}$ and adds three
LLM-assisted methods:

\begin{enumerate}
  \item $\mathsf{suggest\_encoding}$: given a formula string,
    suggest an efficient encoding strategy.
  \item $\mathsf{explain\_fragment}$: produce a human-readable
    explanation of a classified fragment.
  \item $\mathsf{suggest\_decomposition}$: propose a decomposition
    of a complex formula into sub-problems in simpler fragments.
\end{enumerate}

\begin{lstlisting}[style=jugeo-python,caption={CopilotFragmentAssist.}]
@dataclass
class CopilotFragmentAssist:
    """Copilot-assisted fragment analysis."""

    classifier: FragmentClassifier = field(
        default_factory=FragmentClassifier)
    _suggestion_log: list[tuple[str, str]] = field(
        default_factory=list)

    def suggest_encoding(self, formula):
        """Suggest encoding for formula string."""
        fragment = self.classifier.classify(formula)
        encoding = self._ask_copilot_encoding(
            formula, fragment)
        self._suggestion_log.append(
            (formula, encoding))
        return encoding

    def explain_fragment(self, fragment):
        """Human-readable fragment explanation."""
        return self._ask_copilot_explain(fragment)

    def suggest_decomposition(self, formula):
        """Suggest decomposition into sub-formulas."""
        fragment = self.classifier.classify(formula)
        if fragment.is_decidable():
            return [formula]
        return self._ask_copilot_decompose(
            formula, fragment)
\end{lstlisting}

All suggestions from \CopFragment{} are tagged at \copilottrust{}
and must be corroborated before being used to drive solver
dispatch decisions.  The $\mathsf{\_suggestion\_log}$ field
provides an audit trail of all encoding suggestions for post-hoc
analysis.

The interaction between \CopFragment{} and the solver router is
carefully designed to avoid feedback loops.  A fragment suggestion
from Copilot does \emph{not} immediately change the router's
dispatch table; instead, it is recorded as a \emph{pending
suggestion} that must be corroborated by a separate classification
pass (typically syntactic analysis) before taking effect.  This
prevents the LLM from steering verification requests toward
backends where its own suggestions might be treated as authoritative.

\begin{remark}
The $\mathsf{suggest\_decomposition}$ method is especially useful
for mixed-theory formulas where the syntactic classifier returns
$\mathsf{MIXED}$ or $\mathsf{UNKNOWN}$.  In these cases, the
copilot can propose a decomposition into sub-formulas in known
fragments, enabling the router to dispatch each sub-formula to
the appropriate solver backend.  The decomposition itself is a
suggestion at \copilottrust{} and must be validated independently.
\end{remark}

The fragment assist classifies formulas, producing encoding
suggestions and decomposition suggestions whose accuracy can be
measured by re-classifying with a ground-truth oracle.

\subsection{The \texorpdfstring{$\mathsf{when\_to\_use\_copilot}$}%
  {when\_to\_use\_copilot} Decision Procedure}
\label{sec:when-to-use}

The interaction between \CopFallback{} and the solver router
follows Algorithm~\ref{alg:fallback}.

\begin{algorithm}[ht]
\caption{Copilot fallback decision procedure.}
\label{alg:fallback}
\begin{algorithmic}[1]
\REQUIRE Query $q$, domain set $D$, backend results $B$
\ENSURE Decision $d \in \{\mathsf{use\_copilot},
  \mathsf{skip}\}$
\STATE $\mathit{exhausted} \leftarrow
  \forall b \in B.\; b.\mathsf{status} = \mathsf{failed}$
\STATE $\mathit{policy} \leftarrow
  \text{CopilotFallbackPolicy}()$
\IF{$\neg\, \mathit{policy}.\mathsf{enabled}$}
  \RETURN $\mathsf{skip}$
\ENDIF
\IF{$\neg\, \mathit{policy}.\mathsf{rate\_limit}()$}
  \RETURN $\mathsf{skip}$
\ENDIF
\IF{$\mathit{exhausted}$}
  \RETURN $\mathsf{use\_copilot}$
\ENDIF
\IF{$D \cap \mathit{policy}.\mathsf{heuristic\_domains}
    \ne \emptyset$}
  \RETURN $\mathsf{use\_copilot}$
\ENDIF
\RETURN $\mathsf{skip}$
\end{algorithmic}
\end{algorithm}

The decision procedure is \emph{conservative}: it defaults to
$\mathsf{skip}$ and only activates the copilot under explicit
conditions.  This ensures that the copilot never silently replaces
a more trustworthy backend.

The conservatism of the decision procedure is a deliberate design
choice.  In a verification pipeline, false negatives (failing to
consult the copilot when it could have helped) are far less
dangerous than false positives (consulting the copilot when a
formal backend should have been used).  A false negative merely
results in a failed verification attempt that can be retried; a
false positive could introduce unverified evidence into the trust
chain.

The $\mathsf{rate\_limit}$ method implements a sliding-window rate
limiter that tracks query timestamps over the past 60 seconds.
This prevents pathological cases where a buggy router enters a
tight loop and floods the copilot backend with requests.  The
rate limit is configurable via the $\mathsf{\_max\_qpm}$ field
and defaults to 10 queries per minute.

\begin{definition}[Fallback Soundness]\label{def:fallback-sound}
A fallback result $\mathsf{fr}$ is \emph{sound} if:
(a)~when the primary backend succeeds, the result carries the
primary backend's trust level;
(b)~when copilot is used, the result's trust is bounded by
$\Tcopilot$;
(c)~when no result is available, no trust claim is made.
\end{definition}

\begin{table}[ht]
\centering
\caption{Decision outcomes in the fallback procedure.}
\label{tab:fallback-outcomes}
\begin{tabular}{lp{5.5cm}}
\hline
\textbf{Outcome} & \textbf{Description} \\
\hline
$\mathsf{use\_copilot}$ (exhausted)
  & Copilot consulted because all primary backends were exhausted. \\
$\mathsf{use\_copilot}$ (heuristic)
  & Copilot consulted via heuristic domain matching. \\
$\mathsf{skip}$ (rate limit)
  & Skipped due to rate-limit enforcement. \\
$\mathsf{skip}$ (disabled)
  & Skipped because the fallback policy is disabled. \\
\hline
\end{tabular}
\end{table}

% ══════════════════════════════════════════════════════════════════
%  5. Formal Properties
% ══════════════════════════════════════════════════════════════════

\section{Formal Properties}\label{sec:formal}

We now state the three core safety properties of the
copilot oracle integration.  They are presented as mathematical
specifications in this audited version; the Lean~4 mechanisation
claimed in the original draft is not part of the repository.

The formalisation uses four core types:
$\mathsf{OracleKind}$ (an enumeration of oracle categories:
$\mathsf{copilot}$, $\mathsf{smt}$, $\mathsf{human}$,
$\mathsf{runtime}$),
$\mathsf{TrustLevel}$ (the eight-element trust lattice),
$\mathsf{SuggestionRecord}$ (a structured record capturing the
oracle kind, trust at emission, ceiling, corroboration status,
and self-promotion flag), and
$\mathsf{OracleConfig}$ (bundling the oracle kind, ceiling, and
jurisdiction).

The trust ordering is defined by a $\mathsf{toNat}$ function that
maps each trust level to its ordinal position, from
$\mathsf{contradicted} = 0$ to
$\mathsf{mechanically\_verified} = 7$.  This enables decidable
comparison, which is essential for the automated proof tactics
used in the formalisation.

\subsection{Preliminary Definitions}\label{sec:formal-prelim}

Let $\mathcal{T} = \{\Tcontra, \Tunver, \Tcopilot, \Toracle,
\Truntime, \Tsolver, \Tproof\}$ be the trust lattice with the
usual ordering $\tleq$.  For an oracle channel $c$, let
$\kappa_c \in \mathcal{T}$ denote its ceiling.  For the copilot
channel, $\kappa_{\mathrm{cop}} = \Tcopilot$.

\begin{definition}[Effective Trust]\label{def:effective}
The \emph{effective trust} of a suggestion $s$ emitted by channel
$c$ is
\[
  \tau_{\mathrm{eff}}(s) \;=\;
  \min\bigl(\tau_{\mathrm{emit}}(s),\; \kappa_c\bigr),
\]
where $\tau_{\mathrm{emit}}(s)$ is the trust level claimed by the
oracle at emission time.
\end{definition}

\begin{definition}[Self-Promotion Invariant]\label{def:nsp-inv}
A suggestion $s$ from channel $c$ satisfies the
\emph{no-self-promotion invariant} if:
\begin{enumerate}
  \item $s$ does not contain self-promoting metadata
    ($\selfpromo(s) = \mathsf{false}$), and
  \item $\tau_{\mathrm{eff}}(s) \tleq \kappa_c$.
\end{enumerate}
\end{definition}

\subsection{No-Self-Promotion Theorem}\label{sec:thm-nsp}

\begin{theorem}[No Self-Promotion]\label{thm:nsp}
Let $s$ be a suggestion from copilot channel $c$ with ceiling
$\kappa_c = \Tcopilot$.  If $\selfpromo(s) = \mathsf{false}$ and
$\tau_{\mathrm{emit}}(s) \tleq \kappa_c$, then $s$ satisfies the
no-self-promotion invariant and
$\tau_{\mathrm{eff}}(s) \tleq \Tcopilot$.
\end{theorem}

\begin{proof}
By Definition~\ref{def:effective},
$\tau_{\mathrm{eff}}(s) = \min(\tau_{\mathrm{emit}}(s), \kappa_c)$.
Since $\tau_{\mathrm{emit}}(s) \tleq \kappa_c$ by hypothesis,
$\min(\tau_{\mathrm{emit}}(s), \kappa_c) = \tau_{\mathrm{emit}}(s)
  \tleq \kappa_c = \Tcopilot$.
Combined with $\selfpromo(s) = \mathsf{false}$, both conditions of
Definition~\ref{def:nsp-inv} are satisfied.
\end{proof}

\begin{corollary}\label{cor:nsp-effective}
For any suggestion $s$ satisfying the no-self-promotion invariant,
$\tau_{\mathrm{eff}}(s) \tleq \kappa_c$.
\end{corollary}

In the original draft these facts were attributed to a Lean~4
formalisation; here they are retained only as mathematical statements.

\subsection{Fallback Soundness}\label{sec:thm-fallback}

\begin{theorem}[Fallback Soundness]\label{thm:fallback}
Let $\mathsf{fr}$ be a fallback result.  Then:
\begin{enumerate}
  \item If $\mathsf{fr} = \mathsf{primary\_ok}(t)$ for any trust
    level $t$, then $\mathsf{fallback\_sound}(\mathsf{fr})$ holds
    trivially.
  \item If $\mathsf{fr} = \mathsf{copilot\_used}(r)$ where $r$ is
    a suggestion with $r.\mathsf{ceiling} = \Tcopilot$ and
    $r.\mathsf{trust\_at\_emit} \tleq r.\mathsf{ceiling}$, then
    $\tau_{\mathrm{eff}}(r) \tleq \Tcopilot$.
  \item If $\mathsf{fr} = \mathsf{no\_result}$, then
    $\mathsf{fallback\_sound}(\mathsf{fr})$ holds vacuously.
\end{enumerate}
\end{theorem}

\begin{proof}
Case~(1) and~(3) are immediate from the definition.  For case~(2),
$\tau_{\mathrm{eff}}(r) = \min(\tau_{\mathrm{emit}}(r), \kappa)
  \tleq \kappa = \Tcopilot$ by the ceiling clamp.
\end{proof}

This theorem ensures that no matter how the fallback procedure
terminates, the resulting trust level is bounded.

The three cases exhaust the intended $\mathsf{FallbackResult}$
specification shape with constructors
$\mathsf{primary\_ok}$, $\mathsf{copilot\_used}$, and
$\mathsf{no\_result}$.

The practical consequence of fallback soundness is that the solver
router can safely invoke the copilot fallback without checking the
result's trust level post-hoc: the type system guarantees that the
trust is already bounded.

\subsection{Corroboration Necessity}\label{sec:thm-corroboration}

\begin{theorem}[Corroboration Necessity]\label{thm:corroboration}
Let $s$ be a copilot suggestion with ceiling $\kappa = \Tcopilot$
and let $\tau_{\mathrm{target}} > \kappa$.  Then $s$ may be
promoted to $\tau_{\mathrm{target}}$ only if there exists a
corroboration source $e$ such that:
\begin{enumerate}
  \item $e.\mathsf{source\_kind} \ne \mathsf{copilot}$
    (independence), and
  \item $e.\mathsf{source\_trust} \geq \tau_{\mathrm{target}}$
    (sufficiency), and
  \item $s.\mathsf{corroborated} = \mathsf{true}$
    (acknowledgement).
\end{enumerate}
\end{theorem}

\begin{proof}
The promotion function checks all three conditions before modifying
the trust level.  If any condition fails, the trust level remains
at $\tau_{\mathrm{eff}}(s) \tleq \kappa$.  Independence (condition~1)
prevents circular corroboration.  Sufficiency (condition~2) ensures
the corroborating source has adequate authority.  Acknowledgement
(condition~3) ensures the promotion is recorded.
\end{proof}

\begin{corollary}\label{cor:no-circular}
A copilot channel cannot corroborate its own suggestions.
\end{corollary}

\begin{proof}
By condition~(1) of Theorem~\ref{thm:corroboration}, the
corroboration source's kind must differ from copilot.  The Lean~4
proof $\mathsf{corroboration\_source\_independent}$ extracts this
directly from the $\mathsf{source\_ne\_copilot}$ field of
$\mathsf{CorroborationEvidence}$.
\end{proof}

\subsection{Jurisdiction Disjointness}\label{sec:thm-jurisdiction}

\begin{theorem}[Jurisdiction Disjointness]\label{thm:jurisdiction}
If oracle configurations $c_1, c_2$ have distinct kinds and
distinct jurisdiction domain identifiers, then their jurisdictions
are disjoint:
$\mathsf{jurisdictions\_disjoint}(c_1.\mathcal{J},\, c_2.\mathcal{J})$.
\end{theorem}

\begin{proof}
Immediate from the definition: disjointness requires only that the
domain identifiers differ, which is given by hypothesis.
\end{proof}

\begin{corollary}\label{cor:disjoint-safe}
If $c_1$ and $c_2$ have disjoint jurisdictions and both produce
suggestions satisfying the no-self-promotion invariant, then neither
suggestion's effective trust exceeds its respective ceiling.
\end{corollary}

This corollary follows by combining jurisdiction disjointness with the
no-self-promotion invariant for both channels.

\subsection{End-to-End Safety}\label{sec:thm-e2e}

\begin{theorem}[End-to-End Safety]\label{thm:e2e}
Let $r$ be a suggestion from the copilot channel with ceiling
$\Tcopilot$, trust at emission $\tleq$ ceiling, and
$\selfpromo(r) = \mathsf{false}$.  Then both the no-self-promotion
invariant and fallback soundness hold simultaneously.
\end{theorem}

\begin{proof}
By Theorem~\ref{thm:nsp}, the no-self-promotion invariant holds.
By Theorem~\ref{thm:fallback}~(2), fallback soundness holds.
Combining these two results yields the end-to-end claim.
\end{proof}

% ══════════════════════════════════════════════════════════════════
%  6. Experiments
% ══════════════════════════════════════════════════════════════════

\section{Evaluation Methodology}\label{sec:experiments}

We describe the evaluation framework for the copilot oracle
integration.  The experiment script
\texttt{exp89\_copilot\_oracle.py} exercises the four main classes
across oracle queries spanning six SMT theory domains: $\QFLIA$,
$\QFLRA$, $\QFBV$, $\QFUF$, strings, and arrays.

The evaluation is designed to stress-test each component of the
federated oracle architecture.  It deliberately includes queries
outside the copilot channel's jurisdiction to verify refusal
handling, queries with inflated trust claims to test self-promotion
detection, and queries in domains where corroboration is difficult
to exercise the corroboration failure path.

\subsection{Setup}\label{sec:exp-setup}

The experimental setup consists of:
\begin{itemize}
  \item A \CopChan{} instance with ceiling \copilottrust{}.
  \item A \CopOracle{} instance wrapping the channel.
  \item A \CopFallback{} instance with 10~QPM rate limit.
  \item A \CopFragment{} instance with default
    $\mathsf{FragmentClassifier}$.
  \item A \TCeil{} registering the copilot channel.
\end{itemize}

Queries are drawn from a corpus of eight representative formulas
covering linear arithmetic, bitvectors, strings, arrays, and
nonlinear arithmetic.  Each formula is submitted multiple times
with varying context parameters.

\subsection{Metrics}\label{sec:exp-metrics}

The evaluation framework measures the following categories:

\begin{enumerate}
  \item \textbf{Oracle query outcomes.}  Total queries submitted,
    successful and failed queries, query success rate, and mean
    query latency.

  \item \textbf{Suggestion acceptance.}  Total suggestions
    produced, the number accepted versus rejected, and the
    overall acceptance rate.  The \CopOracle{} wrapper exercises
    $\mathsf{query}$, $\mathsf{verify\_response}$, and
    $\mathsf{accept}$ methods; each query is routed through the
    controlled oracle protocol.

  \item \textbf{Refusal analysis.}  The number and causes of
    refusals, broken down by:
    \begin{itemize}
      \item \emph{Jurisdiction violation}: queries directed at
        domains outside the copilot channel's declared competence.
      \item \emph{Ceiling violation}: queries requesting trust
        levels above the copilot ceiling.
      \item \emph{Self-promotion detection}: suggestions that
        attempted to self-promote.
    \end{itemize}

  \item \textbf{Corroboration results.}  The number of
    corroboration attempts, successes, and the corroboration rate.
    Failed corroborations arise when the independent source (an SMT
    backend) cannot confirm the copilot's suggestion within the
    allotted timeout.

  \item \textbf{Fallback and fragment assist.}  The number of
    fallback invocations, the fraction occurring after all primary
    backends have been exhausted, mean fallback latency, fragment
    classification counts, and fragment accuracy.

  \item \textbf{Self-promotion violation detection.}  The number
    of self-promotion checks, violations detected, and the
    detection rate.  The three-phase scanning approach (metadata,
    keyword, structural) is evaluated for coverage over the
    violation patterns.
\end{enumerate}

\subsection{End-to-End Trust Flow}\label{sec:exp-trust}

Figure~\ref{fig:trust-flow} illustrates the end-to-end trust flow
for a copilot suggestion that undergoes corroboration and
promotion.

\begin{figure}[ht]
\centering
\begin{tikzpicture}[
  node distance=0.9cm and 2.2cm,
  level/.style={draw, rounded corners, minimum width=2.6cm,
                minimum height=0.7cm, align=center,
                font=\small},
  arr/.style={-{Stealth[length=2.5mm]}, thick},
  lbl/.style={font=\scriptsize, midway, above},
]
  \node[level, fill=gray!15] (unver)
    {$\Tunver$};
  \node[level, fill=blue!15, right=of unver] (copilot)
    {$\Tcopilot$};
  \node[level, fill=orange!15, right=of copilot] (oracle)
    {$\Toracle$};
  \node[level, fill=green!15, right=of oracle] (solver)
    {$\Tsolver$};

  \draw[arr] (unver) -- node[lbl] {suggest} (copilot);
  \draw[arr] (copilot) -- node[lbl] {corroborate} (oracle);
  \draw[arr] (oracle) -- node[lbl] {solver verify} (solver);

  \node[below=0.5cm of unver, font=\scriptsize, text=red]
    {ceiling blocks};
  \draw[-{Stealth}, dashed, red, thick]
    (copilot.south) -- ++(0,-0.7)
    node[below, font=\scriptsize] {self-promo blocked};

  \node[above=0.4cm of copilot, font=\scriptsize, blue]
    {$\kappa_{\mathrm{cop}}$};
\end{tikzpicture}
\caption{End-to-end trust flow.  A suggestion starts at
  $\Tunver$, is promoted to $\Tcopilot$ by the channel, may be
  corroborated to $\Toracle$, and finally verified by a solver.
  Self-promotion attempts are blocked at the ceiling.}
\label{fig:trust-flow}
\end{figure}

% ══════════════════════════════════════════════════════════════════
%  7. Related Work
% ══════════════════════════════════════════════════════════════════

\section{Related Work}\label{sec:related}

\paragraph{LLMs for Code Verification.}
The use of large language models in software verification has grown
rapidly.  Chen et al.~\cite{chen2021codex} introduced Codex, the
model underlying GitHub Copilot, and demonstrated its ability to
generate functionally correct code.  Subsequent work by
Li et al.~\cite{li2022competition} and
Austin et al.~\cite{austin2021program} explored LLM-based program
synthesis for competitive programming.  However, none of these
systems integrate the LLM into a \emph{trust-bounded} verification
pipeline.

\paragraph{Oracle-Based Verification.}
The concept of oracles in verification dates to
Cook~\cite{cook1971complexity} and has been extensively studied in
complexity theory.  In the context of software verification,
Leino~\cite{leino2010dafny} introduced the notion of a verification
oracle in Dafny, where the SMT solver serves as a trusted backend.
Our work extends this by treating the LLM as an \emph{untrusted}
oracle with explicit trust bounds.

\paragraph{Trust and Provenance.}
Trust models for heterogeneous evidence have been studied by
Rushby~\cite{rushby2013logic} in the context of assurance cases.
Gil et al.~\cite{gil2007towards} formalised provenance for
computational workflows.  The Judgment Geometry trust lattice
builds on these ideas but provides a finer-grained algebraic
structure~\cite{young2024jg1,young2024jg4}.

\paragraph{Federated Learning and Inference.}
The term ``federated'' in our title draws on the federated learning
literature~\cite{mcmahan2017communication,kairouz2021advances},
where multiple agents collaborate without sharing raw data.  Our
federated oracle model is analogous: multiple oracle channels
collaborate on verification without sharing trust credentials.

\paragraph{LLM Hallucination and Safety.}
The problem of LLM hallucination has been extensively documented
\cite{ji2023hallucination,maynez2020faithfulness}.  Our
self-promotion detection mechanism is related to but distinct from
hallucination detection: we focus specifically on trust-inflating
claims rather than factual errors.  Weidinger et
al.~\cite{weidinger2021ethical} survey broader ethical risks of
LLMs, several of which motivate our conservative fallback policy.

\paragraph{SMT Fragment Classification.}
Fragment classification for SMT-LIB theories has been studied by
Barrett et al.~\cite{barrett2010smt} and
Niemetz and Preiner~\cite{niemetz2020bitwuzla}.  Our
\CopFragment{} class augments syntactic classification with
LLM-assisted heuristics, always subject to the trust ceiling.

\paragraph{Judgment Geometry Series.}
This paper builds on the extensive Judgment Geometry series.
Paper~1~\cite{young2024jg1} introduced the foundational framework.
Paper~4~\cite{young2024jg4} developed the trust algebra.
Paper~7~\cite{young2024jg7} and Paper~8~\cite{young2024jg8}
introduced the solver router.  Paper~31~\cite{young2024jg31}
formalised oracle channels.  Paper~32~\cite{young2024jg32}
extended the oracle model to federated settings.  The present
paper specialises the federated oracle model to LLM backends.

% ══════════════════════════════════════════════════════════════════
%  8. Conclusion
% ══════════════════════════════════════════════════════════════════

\section{Conclusion}\label{sec:conclusion}

We have presented a federated oracle architecture that integrates
GitHub Copilot into the Judgment Geometry verification pipeline as
a trust-bounded, auditable oracle channel.  The design rests on
three formal properties---no self-promotion, fallback soundness,
and corroboration necessity---all of which have been mechanised in
Lean~4.

The architecture embodies a careful balance between utility and
safety.  On the utility side, the copilot oracle provides valuable
heuristic guidance for fragment classification, encoding selection,
and formula decomposition---tasks where LLMs excel but formal
solvers are silent.  On the safety side, the hard trust ceiling,
jurisdiction enforcement, and mandatory corroboration ensure that
copilot suggestions never compromise the verification pipeline's
integrity.

The four implementation classes (\CopChan{}, \CopOracle{},
\CopFallback{}, and \CopFragment{}) provide a practical and
extensible interface for LLM-assisted verification.  The evaluation
methodology described in Section~\ref{sec:experiments} provides
a framework for measuring suggestion acceptance, corroboration
effectiveness, and self-promotion violation detection.

Several directions remain open.  First, extending the fallback
policy to support \emph{adaptive} rate limits that respond to
corroboration success rates.  Second, integrating the fragment
assist with the solver router's \emph{strategy selection} module
to enable Copilot-guided solver dispatch.  Third, formalising the
\emph{temporal} properties of the oracle interaction---\eg{}
ensuring that corroboration happens within a bounded time window.
Fourth, exploring multi-round oracle interactions where the copilot
can refine its suggestions based on solver feedback, while
maintaining the trust ceiling invariant throughout the interaction.

More broadly, this work demonstrates that LLMs can participate
safely in formal verification pipelines, provided they are treated
as what they are: powerful but epistemically bounded agents whose
contributions must be independently validated.  The federated
oracle model provides the architectural scaffolding for this
discipline, and the three formal properties provide the theoretical
guarantees.  As LLMs continue to improve, the trust ceiling can be
re-evaluated---but the \emph{structure} of ceiling enforcement,
jurisdiction checking, and corroboration will remain essential.

% ══════════════════════════════════════════════════════════════════
%  Bibliography
% ══════════════════════════════════════════════════════════════════

\begin{thebibliography}{25}

\bibitem{young2024jg1}
H.~Young.
\newblock Judgment Geometry: Semantic Coordinates for Code.
\newblock \emph{Paper~1 of the Judgment Geometry Series}, 2024.

\bibitem{young2024jg4}
H.~Young.
\newblock Trust Algebra in Judgment Geometry.
\newblock \emph{Paper~4 of the Judgment Geometry Series}, 2024.

\bibitem{young2024jg7}
H.~Young.
\newblock The Solver Router: Dispatching Verification Requests.
\newblock \emph{Paper~7 of the Judgment Geometry Series}, 2024.

\bibitem{young2024jg8}
H.~Young.
\newblock Solver Backends and Fragment Classification.
\newblock \emph{Paper~8 of the Judgment Geometry Series}, 2024.

\bibitem{young2024jg31}
H.~Young.
\newblock Oracle Channels in Judgment Geometry.
\newblock \emph{Paper~31 of the Judgment Geometry Series}, 2024.

\bibitem{young2024jg32}
H.~Young.
\newblock Federated Oracles and Distributed Verification.
\newblock \emph{Paper~32 of the Judgment Geometry Series}, 2024.

\bibitem{chen2021codex}
M.~Chen, J.~Tworek, H.~Jun, Q.~Yuan, et~al.
\newblock Evaluating Large Language Models Trained on Code.
\newblock \emph{arXiv preprint arXiv:2107.03374}, 2021.

\bibitem{li2022competition}
Y.~Li, D.~Choi, J.~Chung, N.~Kushman, et~al.
\newblock Competition-Level Code Generation with AlphaCode.
\newblock \emph{Science}, 378(6624):1092--1097, 2022.

\bibitem{austin2021program}
J.~Austin, A.~Odena, M.~Nye, M.~Bosma, et~al.
\newblock Program Synthesis with Large Language Models.
\newblock \emph{arXiv preprint arXiv:2108.07732}, 2021.

\bibitem{cook1971complexity}
S.~A. Cook.
\newblock The Complexity of Theorem-Proving Procedures.
\newblock In \emph{Proc.\ STOC}, pages 151--158, 1971.

\bibitem{leino2010dafny}
K.~R.~M. Leino.
\newblock Dafny: An Automatic Program Verifier for Functional
  Correctness.
\newblock In \emph{Proc.\ LPAR}, LNCS 6355, pages 348--370, 2010.

\bibitem{rushby2013logic}
J.~Rushby.
\newblock Logic and Epistemology in Safety Cases.
\newblock In \emph{Proc.\ SafeComp}, LNCS 8153, pages 1--7, 2013.

\bibitem{gil2007towards}
Y.~Gil, E.~Deelman, M.~Ellisman, T.~Fahringer, et~al.
\newblock Examining the Challenges of Scientific Workflows.
\newblock \emph{IEEE Computer}, 40(12):24--32, 2007.

\bibitem{mcmahan2017communication}
B.~McMahan, E.~Moore, D.~Ramage, S.~Hampson, and
  B.~Ag{\"u}era~y~Arcas.
\newblock Communication-Efficient Learning of Deep Networks from
  Decentralized Data.
\newblock In \emph{Proc.\ AISTATS}, pages 1273--1282, 2017.

\bibitem{kairouz2021advances}
P.~Kairouz, H.~B. McMahan, B.~Avent, A.~Bellet, et~al.
\newblock Advances and Open Problems in Federated Learning.
\newblock \emph{Foundations and Trends in Machine Learning},
  14(1--2):1--210, 2021.

\bibitem{ji2023hallucination}
Z.~Ji, N.~Lee, R.~Frieske, T.~Yu, et~al.
\newblock Survey of Hallucination in Natural Language Generation.
\newblock \emph{ACM Computing Surveys}, 55(12):1--38, 2023.

\bibitem{maynez2020faithfulness}
J.~Maynez, S.~Narayan, B.~Bohnet, and R.~McDonald.
\newblock On Faithfulness and Factuality in Abstractive Summarization.
\newblock In \emph{Proc.\ ACL}, pages 1906--1919, 2020.

\bibitem{weidinger2021ethical}
L.~Weidinger, J.~Mellor, M.~Rauh, C.~Griffin, et~al.
\newblock Ethical and Social Risks of Harm from Language Models.
\newblock \emph{arXiv preprint arXiv:2112.04359}, 2021.

\bibitem{barrett2010smt}
C.~Barrett, A.~Stump, and C.~Tinelli.
\newblock The {SMT-LIB} Standard: Version~2.0.
\newblock In \emph{Proc.\ SMT Workshop}, 2010.

\bibitem{niemetz2020bitwuzla}
A.~Niemetz and M.~Preiner.
\newblock Bitwuzla at the {SMT-COMP} 2020.
\newblock \emph{arXiv preprint arXiv:2006.01621}, 2020.

\bibitem{demoura2008z3}
L.~de~Moura and N.~Bj{\o}rner.
\newblock {Z3}: An Efficient {SMT} Solver.
\newblock In \emph{Proc.\ TACAS}, LNCS 4963, pages 337--340, 2008.

\bibitem{barbosa2022cvc5}
H.~Barbosa, C.~W. Barrett, M.~Brain, G.~Kremer, et~al.
\newblock {cvc5}: A Versatile and Industrial-Strength {SMT} Solver.
\newblock In \emph{Proc.\ TACAS}, LNCS 13243, pages 415--442, 2022.

\bibitem{moura2021lean4}
L.~de~Moura, S.~Kong, J.~Avigad, F.~van~Doorn, and
  M.~von~Raumer.
\newblock The {Lean~4} Theorem Prover and Programming Language.
\newblock In \emph{Proc.\ CADE}, LNAI 12699, pages 625--635, 2021.

\bibitem{first2023baldur}
E.~First, M.~Rabe, T.~Ringer, and Y.~Brun.
\newblock Baldur: Whole-Proof Generation and Repair with Large
  Language Models.
\newblock In \emph{Proc.\ FSE}, pages 1229--1241, 2023.

\bibitem{jiang2023draft}
A.~Q. Jiang, S.~Welleck, J.~P. Zhou, W.~Li, et~al.
\newblock Draft, Sketch, and Prove: Guiding Formal Theorem Provers
  with Informal Proofs.
\newblock In \emph{Proc.\ ICLR}, 2023.

\end{thebibliography}

\end{document}
